% Figure for Classical Mechanics §A1.5 (left: scalar product as projection; right: vector product perpendicular to the parallelogram spanned by A and B).
\begin{tikzpicture}[>={Stealth[length=2.4mm]}, font=\small]
  % ---- left panel: projection
  \coordinate (O) at (0,0);
  \coordinate (A) at (3.6,0);
  \coordinate (B) at ({2.6*cos(50)},{2.6*sin(50)});
  \coordinate (P) at ({2.6*cos(50)},0);
  \draw[dashed, black!55] (B) -- (P);
  \draw[black!55] ($(P)+(0,0.18)$) -- ++(-0.18,0) -- ++(0,-0.18);
  \draw[->, very thick, figB] (O) -- (A) node[right] {$\vec A$};
  \draw[->, very thick, figB] (O) -- (B) node[above] {$\vec B$};
  \draw[black!55, thin] (0.7,0) arc[start angle=0, end angle=50, radius=0.7];
  \node[font=\scriptsize] at (0.95,0.38) {$\varphi$};
  \draw[|-|, very thick, figR] (0,-0.22) -- ($(P)+(0,-0.22)$) node[midway, below, font=\scriptsize] {$B\cos\varphi$};
  \fill (O) circle (1.8pt);
  % ---- right panel: cross product (oblique view; A and B in the horizontal plane)
  \begin{scope}[xshift=6.6cm, yshift=-0.6cm]
    \coordinate (Q) at (0,0);
    \coordinate (a) at (2.8,-0.5);
    \coordinate (b) at (1.3,1.0);
    \fill[figB!12] (Q) -- (a) -- ($(a)+(b)$) -- (b) -- cycle;
    \draw[dashed, black!45] (a) -- ($(a)+(b)$) -- (b);
    \draw[->, very thick, figB] (Q) -- (a) node[below right] {$\vec A$};
    \draw[->, very thick, figB] (Q) -- (b) node[above left, xshift=2pt] {$\vec B$};
    \draw[->, thin, black] (0.9,-0.16) to[bend right=35] (0.55,0.42);
    \node[font=\scriptsize] at (1.15,0.25) {$\varphi$};
    \draw[->, very thick, figR] (Q) -- (0,2.9) node[above] {$\vec A \times \vec B$};
    \draw[black!55] (0,0.25) -- (0.14,0.22) -- (0.14,-0.03);
    \node[font=\scriptsize, black!70, align=center] at (2.3,0.25) {area\\ $AB\sin\varphi$};
    \fill (Q) circle (1.8pt);
  \end{scope}
\end{tikzpicture}
